\documentclass[11pt]{article}
\usepackage[margin=25mm]{geometry}
\usepackage{fontspec}
\setmainfont{Latin Modern Roman}
\usepackage{amsmath,amssymb,mathtools}
\usepackage{unicode-math}
\setmathfont{Latin Modern Math}
\usepackage{xurl,hyperref}
\hypersetup{colorlinks=true,urlcolor=blue}
\setcounter{secnumdepth}{0}
\setlength{\parindent}{0pt}
\setlength{\parskip}{5pt}
\setlength{\emergencystretch}{3em}
\newcommand{\sref}[2]{\hyperlink{source:#1:#2}{[S#2]}}
\newcommand{\cataloguescope}[1]{\paragraph{Recorded target.}#1}
\newcommand{\R}{\mathbb R}
\newcommand{\C}{\mathbb C}
\newcommand{\N}{\mathbb N}
\newcommand{\Z}{\mathbb Z}
\newcommand{\Q}{\mathbb Q}
\newcommand{\F}{\mathbb F}
\newcommand{\eps}{\varepsilon}
\begin{document}
\section{Fluctuation Limits for Moderately Interacting Particle Systems}\label{statement}

\subsection{Definitions and mathematical statement}

\paragraph{Shared notation.}$\mathbb N=\{1,2,\ldots\}$; the standard number fields have their usual meanings. All additional conventions are stated locally.


For particle positions $X_i(t)$, the empirical measure is $\mu_N(t)=N^{-1}\sum_{i=1}^N\delta_{X_i(t)}$. The source uses $V_\eta=Z_\eta*Z_\eta$, a smooth normalized symmetric compactly supported convolution-square potential, with algebraic interaction length $\eta=N^{-\beta}$. Its intermediate nonlocal density $u^\eta$ solves the report's aggregation--diffusion equation. Smoothed measures are $f^{N,\eta}=\mu^{N,\eta}*Z_\eta$ and $g^\eta=u^\eta*Z_\eta$. Intermediate fluctuations are $\sqrt N(\mu^{N,\eta}-u^\eta)$, distinct from fluctuations around the final local limit $u$.

\paragraph{Catalogue wording retained for traceability.}
Prove mean-field convergence for Keller–Segel-type and Dirac-approximating potentials, and obtain the better-than-$N^{-1/2}$ mean-square rate needed to derive intermediate fluctuation limits from smoothed empirical measures.

\paragraph{Scope and source qualification.}The report already proves a better-than-$N^{-1/2}$ smoothed error estimate conditional on an assumed mean-field convergence result. The remaining task is to verify that premise for the specified additional potentials and derive the fluctuation limit. No interchangeable interpretation of mean-square error and its square root is assumed.
\subsection{Short English statement}
Obtain quantitative mean-field limits strong enough to control intermediate fluctuations for the stated singular or shrinking-range interaction potentials.
\subsection{Research attempt}
\paragraph{Model, date and investigated scope.}
OpenAI GPT-6 Astra (configured model \texttt{gpt-6-astra}), ultra reasoning; generated 11 October 2026. Derived the exact independent smoothed empirical-measure mean-square rate and its bandwidth threshold; no interacting mean-field premise or fluctuation limit.

\paragraph{Sources and literature check.}
Read the supplied definition, catalogue wording, source qualifications and reference entries for this question. The historical source/context and inherited status are retained; no fresh literature or novelty certification is claimed. No external theorem is used without proof in this attempt; explicit model assumptions are distinguished from facts about the source model.

\paragraph{The independent sampling baseline.} Let \(X_1,\ldots,X_N\) be independent samples of a probability density \(\rho\), and put \(Z_\eta(x)=\eta^{-d}Z(x/\eta)\) for a normalized nonnegative \(L^2\) smoothing kernel. With \(f_N=N^{-1}\sum_iZ_\eta(\cdot-X_i)\) and \(g=\rho*Z_\eta\), expansion of the squared norm and independence give
\[
 \mathbb E\|f_N-g\|_2^2
 =\frac1N\left(\|Z_\eta\|_2^2-\|g\|_2^2\right)
 =\frac1N\left(\eta^{-d}\|Z\|_2^2-\|\rho*Z_\eta\|_2^2\right).
\]
The cross terms for distinct particles vanish after centering, and translation leaves each kernel norm unchanged. If \(\rho\in L^2\), Jensen's inequality for normalized nonnegative \(Z_\eta\) bounds \(\|\rho*Z_\eta\|_2\le\|\rho\|_2\). For \(\eta=N^{-\beta}\), \(\beta>0\), the leading mean-square scale is therefore \(N^{-1+d\beta}\).
\paragraph{Rate audit.} This baseline is better than \(N^{-1/2}\) precisely when \(d\beta<1/2\); its square root has a different exponent. It is a sampling calculation, not a bound on a coupled interacting-particle error. Centering at the nonlocal density \(u^\eta\) rather than the final local limit also avoids an additional deterministic bias only if that intermediate evolution is actually controlled.
\paragraph{Source gap.} The report's improved estimate is already conditional on mean-field convergence. The first missing step here is establishing that premise for the singular/shrinking-range potentials with enough uniformity, then proving fluctuation tightness and identifying the limit. Independence was used exactly where the interacting system creates the hard covariance terms.

\paragraph{Verification and adversarial audit.}
Checked centering, vanishing cross terms, kernel scaling and mean-square versus root-mean-square exponents; stated the extra positivity assumption for Jensen. This is a GPT-6 Astra Ultra self-audit, not independent human verification. Prompt A drafting followed by prompt B scope/dependency/boundary self-audit. Checked centering, vanishing cross terms, kernel scaling and mean-square versus root-mean-square exponents; stated the extra positivity assumption for Jensen. The remaining gap is: Control interacting covariances for the source potentials uniformly in the shrinking range and establish the prerequisite mean-field convergence and fluctuation limit. Unresolved status and zero original-target completion are retained.

Coordinator checked the centered independent-sampling variance, nonnegative-kernel convolution bound and squared-error bandwidth exponent. This additional check was model-assisted, not independent human verification.

\paragraph{Tools, process and contributions.}
Python 3 and local file-reading/writing tools for manuscript serialization and structural checks; mathematical calculations were symbolic reasoning, with no numerical experiment or theorem prover. Contribution prompt A was followed by prompt B. GPT-6 Astra with ultra reasoning drafted the research and performed its stated self-audit; any additional coordinator contribution is identified in the verification and audit. A GPT-6-based Codex coordinator assembled the manuscript and checked repository integration; its exact variant and reasoning configuration were not exposed. No human mathematical contribution, independent human proof check or formal proof certification is supplied. Complete derivations and the consulted source references are included here; private reasoning transcripts are not published.

\paragraph{Final status and first remaining gap.}
\textbf{UNRESOLVED}. Control interacting covariances for the source potentials uniformly in the shrinking range and establish the prerequisite mean-field convergence and fluctuation limit.

\section{Completion Estimate}
\noindent\textbf{COMPLETION: 0\%}
The proved calculation is an elementary auxiliary result or a test of a restricted strategy; no new portion of the full source target is established.

\subsection{Sources}
\begin{itemize}
\item[\hypertarget{source:statement:1}{S1}] ULAM.AI and the UnsolvedMath Contributors, UnsolvedMath, catalogue entry. Catalogue curation under CC BY 4.0; definitions and English exposition added locally.
\url{https://huggingface.co/datasets/ulamai/UnsolvedMath}.
\item[\hypertarget{source:statement:2}{S2}] Classical and Quantum Mechanical Models of Many-Particle Systems (2023), Oberwolfach Reports, DOI 10.4171/owr/2023/38, pp. 2179--2246. Original report and mathematical context.
\url{https://doi.org/10.4171/owr/2023/38}.
\end{itemize}
\label{end:statement}
\end{document}
